% =====================================================================
% Introduction + Contribution statement
% =====================================================================
\section{Introduction}
\label{sec:intro}

Parameter-efficient fine-tuning with LoRA~\citep{hu2022lora} has made it cheap
to specialize a base model for a task: practitioners maintain libraries of
adapters layered over a frozen foundation checkpoint. But base models are not
static. When a vendor ships a new version of a family --- Gemma~3 to Gemma~4,
Qwen2 to Qwen2.5 --- every adapter trained on the old base must, in current
practice, be retrained from scratch on the original task data. That data is
often private, expensive to re-collect, or simply gone, and the retraining cost
scales with the number of adapters an organization maintains.

We ask whether a fine-tuning delta can instead be treated as a \emph{portable
object}: transplanted onto a new base version without retraining on task data.
The obstacle is that a LoRA adapter is entangled with the base model's
representation space (\S\ref{sec:problem}) --- it reads and writes in the old
model's activation geometry, which the new version does not preserve. Existing
weight-space transfer methods do not address this: task arithmetic and task
vectors~\citep{ilharco2023taskvectors}, TIES-Merging~\citep{yadav2023ties}, and
DARE~\citep{yu2024dare} all assume a \emph{fixed, shared} representation space
across the weights being combined, and break under version drift.

Our framing is that cross-version portability is a \emph{representational
alignment} problem. We measure the drift between base versions, learn a
lightweight linear map between their activation spaces from a small unlabeled
anchor set, and use it to re-express the adapter in the new model's geometry
while preserving its low-rank structure.

\paragraph{Contributions.}
\begin{itemize}
  \item \textbf{A novel problem: cross-version adapter portability.} We
        formalize transplanting a LoRA adapter from base $M_v$ onto a
        same-family successor $M_{v+1}$ using only the two base checkpoints, the
        adapter, and a small unlabeled anchor set --- no task data and no
        retraining (\S\ref{sec:problem}). To our knowledge this problem is
        unaddressed by prior model-merging or task-arithmetic work.
  \item \textbf{An empirical study of representational drift across model
        versions.} Using CKA over anchor activations, we characterize how, where,
        and how much representation geometry shifts between successive base
        versions, and show this drift predicts the failure of naive adapter
        transfer (\S\ref{sec:drift}).
  \item \textbf{A simple, closed-form alignment method.} We derive a per-module
        linear transplant $\Delta W' = P^{\mathrm{out}} \Delta W\, P^{\mathrm{in}}$
        (ridge / orthogonal-Procrustes) that preserves rank $r$ and thus adds no
        inference cost, fit from a single anchor-set forward pass through each
        base model (\S\ref{sec:method}).
  \item \textbf{Evaluation across tasks and a version pair.} We evaluate on
        2--3 fine-tuning tasks and 1--2 family version pairs against a naive-transfer
        lower bound and a retrained oracle upper bound, reporting the fraction of
        the oracle--naive gap recovered at a fraction of oracle compute
        (\S\ref{sec:experiments}).
  \item \textbf{Open release.} We release code, anchor sets, and the alignment
        maps to support reproducibility and follow-on work (\S\ref{sec:repro}).
\end{itemize}
